\section{Introdução}
\label{sec:introducao}

A Atividade XIII solicita um plano de crescimento e escalabilidade para o projeto de startup do grupo, considerando o momento atual do negócio e as oportunidades de mercado \cite{pucminas2026xiii}. A orientação organiza o trabalho em quatro etapas: diagnóstico do estágio atual, duas estratégias de crescimento, estratégias de escalabilidade e plano de ação resumido com recursos e indicadores. Este documento segue essa ordem (Seções~\ref{sec:diagnostico} a \ref{sec:plano}) e acrescenta uma seção de viabilidade e coerência (Seção~\ref{sec:coerencia}), em atenção ao título da atividade.

O plano apoia-se nas atividades anteriores: o protótipo e os testes com usuários \cite{grupo2026a6}, os canais prioritários \cite{grupo2026a8}, o levantamento de custos e a precificação \cite{grupo2026a9}, o tipo societário \cite{grupo2026a10} e a estrutura organizacional \cite{grupo2026a11}. Em particular, os valores numéricos de custos, preços e receitas provêm da Atividade IX, que os classifica como estimativas do grupo e hipóteses a validar. Esses valores são reproduzidos aqui com essa mesma ressalva, e a simulação do Apêndice~A permite que o grupo substitua as premissas por dados reais assim que o piloto os produzir.

Esta versão final resulta da comparação entre as versões preliminares do grupo: manteve-se a modelagem econômica baseada nos dados da Atividade IX e incorporaram-se o desenho de um piloto assistido de oito semanas, as definições operacionais de ativação e retenção, as ressalvas sobre o alcance das evidências e as exigências de proteção de dados.

Dois princípios orientam o texto. O primeiro é a distinção entre crescer e escalar: crescer é aumentar clientes ou receita, mesmo à custa de custos proporcionalmente maiores; escalar é aumentar clientes ou receita sem crescimento proporcional dos custos, conforme a orientação. O segundo é a coerência: uma estratégia só é recomendada se for compatível com o estágio do negócio, com a economia unitária, com os canais já definidos e com a capacidade de uma equipe de quatro integrantes.
